\section{Packet estimates and a near-Cauchy bound}\label{p:section}

We first specify the packet class and the estimates that will be used through
all subsequent decompositions.  In particular, the auxiliary square weights
in the estimates below need not be the squares of the packet amplitudes.
Keeping this distinction is necessary when some branches of a decomposition
have been removed.

\subsection{Conventions and external estimates}

We write $\braket{z}=(1+|z|^2)^{1/2}$.
Fix an integer $k\geq1$, and write
\begin{equation}\label{p:exponents}
 q_k=\frac{2(k+2)}k,\qquad p_k=\frac{2(k+1)}k.
\end{equation}
All frequency patches lie in a fixed bounded set in $\R^k$.  The phase is
$x\cdot v-t|v|^2$, so its normal vector is $(2v,1)$.  Frequency sets
$U_1,\ldots,U_{k+1}$ are $\delta$-transverse if
\begin{equation}\label{p:transversality}
 \big|\det((2v_1,1),\ldots,(2v_{k+1},1))\big|\geq\delta
 \quad(v_j\in U_j).
\end{equation}
Changing the sign of time changes the sign of the paraboloid and changes none
of the estimates.  A positive definite quadratic form with fixed upper and
lower eigenvalue bounds is handled by a fixed linear change of variables.
No estimate for an indefinite quadratic form is being invoked.

\begin{convention}[Packets and counting]\label{p:conventions}
A packet of scale $N\geq1$ has a frequency label $v$ in a grid of mesh
$N^{-1}$ and a unit lattice cell $T$ in the coordinates
\begin{equation}\label{p:native}
 X_{N,v}(x,t)=\left(\frac{x-2vt}{N},\frac{t}{N^2}\right).
\end{equation}
The center of its physical cell is denoted by $z_T$, and $d_T(z)$ is the
distance from $X_{N,v}(z)$ to its native cell.  Thus $|T|=N^{k+2}$.
Choose fixed, sufficiently large exponents $P$ and $P'>P$, and put
$w_T(z)=(1+d_T(z))^{-P}$.  A band-limited packet of amplitude $s_T$ is
\[
 f_T(x,t)=e^{i(x\cdot v-t|v|^2)}F_T(X_{N,v}(x,t)),\qquad
 |F_T(X)|\leq s_T(1+\dist(X,T))^{-P'},
\]
where the Fourier transform of $F_T$ is supported in a fixed box.  Alternatively,
a smooth packet has a native profile supported in a fixed enlargement of
its associated cell $T$, with a fixed sufficiently large number of
derivatives bounded by $s_T$.
The Fourier box, the profile bounds, and the exponents are uniform throughout
a packet family.

The phase-space multiplicity is bounded as follows.  On each slab of length
$N^2$ containing packet centers, there are boundedly many packets per frequency
slot and per spatial intercept cell of side $N$, after transporting their
centers to the central time of that slab.  The bound is relative to each slab;
there is no restriction on the number of slabs.  Frequencies may move within
their slots and centers may be recentered.  Moving all native centers by a
fixed lattice displacement changes this multiplicity by at most a fixed
polynomial in the displacement; such multiplicities can also be split into
bounded-multiplicity layers.  In particular,
\begin{equation}\label{p:count}
 \#\{T:d_T(z)\leq a\}\leq C N^k(1+a)^{C_k},\qquad a\geq1.
\end{equation}
The energy parameter of the family is
\begin{equation}\label{p:energy}
 E=N^{k+2}\sum_Ts_T^2.
\end{equation}
It is a parameter of the decomposition, not an assertion that the packets are
orthogonal in spacetime.

A unit-cube count means a count in a fixed lattice, allowing a bounded number
of translates or bounded enlargements.  There is one witness per counted cube;
the square-envelope and field witnesses are specified explicitly when they
need not coincide.  Estimates are first proved for finite families and finite
sets of test cubes.  Absolutely convergent decompositions, or truncations with
arbitrarily small error on those test cubes, then give the stated generality.
Constants may depend on the fixed profile and multiplicity bounds, never on
the size of the configuration.
\end{convention}

We record the exact forms of two external geometric inputs.  They will also
be used after the packet estimates have been proved.

\begin{proposition}[Weighted endpoint multilinear Kakeya]\label{p:kakeya}
Let $d\geq2$.  For $j=1,\ldots,d$, let $\mathcal T_j$ be a family of infinite
cylinders of radius $h$ in $\R^d$, with unit direction vectors $e(T)$.  Suppose
\[
 |\det(e(T_1),\ldots,e(T_d))|\geq\vartheta>0
 \quad(T_j\in\mathcal T_j).
\]
For nonnegative weights $c_j(T)$,
\begin{equation}\label{p:kakeya-equation}
 \int_{\R^d}\prod_{j=1}^d
 \left(\sum_{T\in\mathcal T_j}c_j(T)\mathbf1_T\right)^{1/(d-1)}
 \leq C_d\vartheta^{-1/(d-1)}h^d
 \prod_{j=1}^d\left(\sum_{T\in\mathcal T_j}c_j(T)\right)^{1/(d-1)}.
\end{equation}
For lattice cubes of side $a\le h$, the same right-hand side bounds
$a^d$ times the corresponding sum of cube suprema, after a fixed
dimensional enlargement of every cylinder.
\end{proposition}
\begin{proof}
The unweighted unit-radius assertion with arbitrary transverse directions is
\citet[Theorem~2]{Guth2010}; the near-coordinate-axis multilinear formulation
originates in \citet{BCT2006}.  Integral weights are obtained by repeating
cylinders.  Dividing by a common denominator gives rational weights, and
monotone approximation gives arbitrary nonnegative weights.  Dilation gives
$h^d$. Every side-$a$ cube meeting a cylinder lies in its enlargement
of radius $h+C_da\lesssim h$. Integration over each cube, of volume
$a^d$, proves the lattice version, also with separate suprema for the
factors. For paraboloid normals in a bounded patch,
normalization to unit vectors changes the determinant by fixed factors, so
\eqref{p:transversality} gives an inverse power of $\delta$.
\end{proof}

For a set $A$, write $\mathcal N_\Delta(A)$ for its covering number by balls of
radius $\Delta$.  This notation is interchangeable with the number of occupied
$\Delta$-cubes up to dimensional constants.

The projection input is He's higher-rank extension of Bourgain's
discretized projection theorem \citep{BourgainProjection2010,He2020}.
Its conclusion is simultaneous for every sufficiently large subset,
so the successful position set may depend on the direction.

\begin{proposition}[Robust hyperplane projection input]\label{p:projection}
Fix $d\geq2$, $0<\alpha<d$, and $\rho>0$.  There are positive constants
$\eta_0,\gamma$ and $\Delta_0$ with the following consequence.  If
$0<\eta\leq\eta_0$, $0<\Delta<\Delta_0$, $A$ lies in a fixed unit ball,
\[
 \Delta^{-\alpha+\eta}\leq\mathcal N_\Delta(A)
 \leq\Delta^{-\alpha-\eta},\qquad
 \mathcal N_\Delta(A\cap B(x,r))
 \leq\Delta^{-\eta}r^\rho\mathcal N_\Delta(A)
 \quad(\Delta\leq r\leq1),
\]
and $\nu$ is a probability law on unit kernel directions satisfying
\begin{equation}\label{p:projection-angular}
 \nu\{v:|v\cdot w|<r\}\leq\Delta^{-\eta}r^\rho
 \quad(|w|=1,\ \Delta\leq r\leq1),
\end{equation}
then some direction in the support of $\nu$ has
\begin{equation}\label{p:projection-conclusion}
 \mathcal N_\Delta(\pi_{v^\perp}A')
 \geq\Delta^{-(d-1)\alpha/d-\gamma}
 \quad\text{for every }A'\subset A\text{ with }
 \mathcal N_\Delta(A')\geq\Delta^\eta\mathcal N_\Delta(A).
\end{equation}
Fixed multiplicative constants and bounded affine normalizations in the
hypotheses are allowed after reducing $\eta_0,\gamma$.
\end{proposition}
\begin{proof}
This is \citet[Theorem~1]{He2020}, with $m=d-1$, after reducing its error
parameter and, if necessary, replacing $\rho$ by a smaller positive number.
Here is the match to its Grassmannian hypothesis.  The competing subspaces
have dimension $d-m=1$; if one is spanned by a unit vector $w$, its Schubert
angle to the projection plane $v^\perp$ is $|v\cdot w|$.  Thus
\eqref{p:projection-angular} is exactly the required nonconcentration near
Schubert cycles.  The theorem excludes a direction set of small $\nu$-mass and
is simultaneous for all sufficiently large subsets of $A$.  Its size and
nonconcentration errors can be made larger than those displayed above, leaving
a fixed positive projection gain $\gamma$.  Shrinking this gain and the error
tolerance absorbs the fixed comparison constants.  Its theorem tolerance is chosen larger than $\eta$: our lower size bound
then implies its required lower bound, and our subset threshold is stronger
than its subset threshold.  No choice of a new size exponent is necessary.
\end{proof}

\subsection{Localization with fixed original regularity}

\begin{lemma}[Reproduction, tails, and unit suprema]\label{p:reproduction}
For the band-limited class in \Cref{p:conventions}, the following operations
are uniform in the size of the packet family.

\begin{enumerate}
\item A polynomially decaying packet is an absolutely convergent sum of packets
with the same fixed enlarged native Fourier support and arbitrarily rapid
Schwartz decay about translated native centers.  The coefficient at a
translation $j\in\Z^{k+1}$ is at most
$C_A s_T(1+|j|)^{-P'}$, where $A$ is the requested output decay order.
\item All fixed native derivative orders satisfy the corresponding envelope
bounds.  If $F$ has Fourier support in a fixed bounded set and $r>0$, then,
for every sufficiently large fixed $L$,
\begin{equation}\label{p:sup-convolution}
 \sup_{|u-z|\leq C}|F(u)|^r
 \leq C_{r,L}\int |F(y)|^r(1+|y-z|)^{-L}\dd y.
\end{equation}
\end{enumerate}
Consequently integral packet estimates below also hold with unit-cube
suprema, or with separate unit suprema for factors of a product.  Output
Schwartz orders may be chosen after a small loss parameter; the original
exponent $P'$ and original smooth-profile derivative order stay fixed.
\end{lemma}
\begin{proof}
Choose a smooth partition $\sum_j\psi_j=1$ on native space, where $\psi_j$ is
supported within fixed distance of $j$.  Choose a compactly supported smooth
Fourier multiplier equal to one on the original Fourier box, and let $K$ be
its Schwartz inverse Fourier transform.  Then
\[
 F_T=\sum_j K*(\psi_jF_T),\qquad
 |\partial^\beta K*(\psi_jF_T)(X)|
 \leq C_{A,\beta}s_T(1+|j|)^{-P'}(1+|X-j|)^{-A}.
\]
These inequalities follow by integrating the Schwartz bounds of $K$ over
the fixed support of $\psi_j$; the decay at $j$ comes from the original
pointwise bound.  They prove absolute convergence and every asserted
output derivative bound.  Convolution without the partition also proves
fixed derivative bounds for the original packet.

For completeness, the maximal inequality is valid also for $r<1$.
Let $M=\sup_y |F(y)|(1+|y-z|)^{-a}$, choosing $a$ so large that $ar>L+k+1$,
and choose a point $y_0$ realizing at least $M/2$.  Reproduction and its
derivative version give
$|\nabla F(y)|\leq C M(1+|y-z|)^a$.
On a fixed sufficiently small ball about $y_0$, this and the choice of
$y_0$ imply $|F(y)|\geq cM(1+|y_0-z|)^a$.
Integrating on that ball proves the weighted maximal estimate with weight
$(1+|y-z|)^{-ar}$, and hence \eqref{p:sup-convolution} after choosing $ar$
large enough.  Truncation of the weighted supremum justifies the argument
when its finiteness has not yet been established.

Translating a native core by $j$ multiplies its weight comparisons and
phase-space multiplicity by a fixed polynomial in $1+|j|$.  For example,
$w_{T+j}(z)\leq C(1+|j|)^Pw_T(z)$.  The sum over $j$ therefore converges
once $P'$ exceeds those fixed powers and $k+1$.  Crucially, $A$ occurs in the
constant $C_A$, not in that required original power $P'$.
For packet estimates use \eqref{p:sup-convolution} and the polynomial
comparison of envelopes at translated points.  For products apply it
separately to each factor and use Tonelli.  The exponent in each convolution
is fixed after the product exponents have been fixed.
\end{proof}

\subsection{Weighted decoupling}

The refined estimate needed here counts where a packet can actually be
active.  In the arbitrary-packet formulations printed in
\citet[Theorem~4.2]{GIOW2020} and
\citet[Theorem~4.2]{DIOWZ2021}, the packet may occupy $2T$ while the
count is made on $T$.  The collar counterexample of
\citet[Theorem~1.1]{SinghSinghParmar2026} rules out that literal
formulation in ambient dimension two.  Their Appendix~A gives a buffered
proof in that dimension.  We give the corresponding calculation in dimension $k+1$, with a
strict support-to-count margin and a common longitudinal slab.  Its only
analytic input is the positive-definite $\ell^2$ decoupling theorem of
\citet{BD2015}.

\begin{lemma}[Buffered refined decoupling]\label{p:buffered-refined}
Put $d=k+1$, $R\geq2$, and $2\leq p\leq q_k$.  For every $\varepsilon>0$
there is $\vartheta_0=\vartheta_0(d,p,\varepsilon)>0$ such that the following
holds for $0<\vartheta\leq\vartheta_0$ and
$\Lambda=R^\vartheta$, with a constant allowed to depend on
$(d,p,\varepsilon,\vartheta)$.  Work in a fixed $O(R)$ test box and one common
padded time slab $S$ of length $O(R)$.  Every tube core meets a fixed
$O(R)$ enlargement of the test box; in particular the family has
polynomial cardinality.  Fine frequency labels lie in a
bounded patch of a translated, bounded-overlap lattice of mesh
$R^{-1/2}$; the allowed class of such lattices is closed under parabolic
rescaling.  A fixed $O_d(1)$ number of cap grids may be separated into
colors.  A labelled tube is the intersection with $S$ of
\[
 \{(x,t):|x-2vt-u|_\infty\leq\Lambda R^{1/2}\}.
\]
For $c\geq1$, $cT$ dilates the transverse width only, in the same common
slab.  Tubes of a fixed label have bounded overlap before clipping by a
spatial box; equivalently, their intercepts are globally separated at
scale $\Lambda R^{1/2}$ up to fixed multiplicity.  Fix
$1\leq a\leq10$ and $a+1/2\leq b\leq20$.  Suppose the nonzero $F_T$
have comparable $L^p$ norms,
\[
 \supp\widehat F_T\subset
 \{(\xi,\tau):|\xi-v|\lesssim R^{-1/2},\
              |\tau+|\xi|^2|\lesssim\Lambda/R\},\qquad
 \|F_T\|_{L^p(\R^d\setminus aT)}\leq R^{-B}\|F_T\|_p,
\]
where $B$ is a sufficiently large fixed number.  Let $U$ be a union of
$R^{1/2}$-cubes in the interior of the box and slab, each meeting at most
$M\geq1$ of the \emph{counted} tubes $bT$.  Then
\begin{equation}\label{p:buffered-equation}
 \left\|\sum_TF_T\right\|_{L^p(U)}
 \leq C_{\varepsilon,\vartheta,d,p}R^\varepsilon
 M^{1/2-1/p}\left(\sum_T\|F_T\|_p^p\right)^{1/p}.
\end{equation}
The same assertion holds for finitely many fixed cap-grid colors.  The
spatial boxes and tube strips may be translated; all overlap assumptions
refer to the unclipped strips.
\end{lemma}
\begin{proof}
We track the dimension-dependent bookkeeping in the induction of
\citet[Appendix~A]{SinghSinghParmar2026}; no refined estimate is used as
an input.  Weighted local Bourgain--Demeter decoupling says that, for any
fixed large $L$, $s\geq1$, $2\leq p\leq q_k$, and a cube $Q_s$ of side $s$,
\begin{equation}\label{p:bd-local}
 \left\|\sum_J H_J\right\|_{L^p(Q_s)}
 \leq D_{\delta,L,d,p}s^\delta
 \left(\sum_J\|H_J\|_{L^p(w_{Q_s,L})}^2\right)^{1/2},\qquad
 w_{Q_s,L}(z)=\left(1+\frac{\dist(z,Q_s)}s\right)^{-L},
\end{equation}
for $s^{-1/2}$ tangential caps with $O(s^{-1})$ normal thickness.  The
weighted version follows by summing the usual local decoupling estimate
over translates of $Q_s$ with rapidly decreasing coefficients.  Its
constant is uniform under the fixed positive-definite affine
normalizations used below.

Let $\mathfrak A_d(R,\Lambda,a,b)$ be the least constant in
\eqref{p:buffered-equation} without $R^\varepsilon$, uniformly for the
stated configurations but with the full domain $1\leq a<b\leq20$.
The theorem starts with $b-a\geq1/2$; smaller positive margins occur
only inside the induction.  Define it for every $\Lambda\geq1$ using the
trivial estimate when $\Lambda\gtrsim R^{1/4}$; it is not defined only
on the curve $\Lambda=R^\vartheta$, since $\Lambda$ stays fixed when $R$
is square-rooted.  For the one-step estimate assume
$R^\vartheta\leq\Lambda\ll R^{1/4}$; this lower bound persists at each
scale reached from the initial $\Lambda=R^\vartheta$.  Group fine labels
into coarse $R^{-1/4}$-cubes $J$, with fixed finite coloring at cap
boundaries.  Translate the midpoint of $S$ to $t=0$, adjusting tube
intercepts accordingly.  If $v_0$ is the center of $J$, assign each fine tube to one
parent box
\[
 \Box_{J,u}=\{(x,t)\in S:
       |x-2v_0t-u|_\infty\leq C_dR^{3/4}\},
 \qquad u\in c_dR^{3/4}\Z^k.
\]
The entire supporting $aT$ lies in its assigned parent box.  This uses
$|v-v_0|=O_d(R^{-1/4})$, the common $O(R)$ time slab, and
$\Lambda R^{1/2}\ll R^{3/4}$.  The boxes have bounded global overlap for
each $J$, and packet families assigned to distinct boxes are disjoint.
On a parent box, remove the carrier and set
\[
 x'=R^{-1/4}(x-2v_0t-u),\qquad t'=R^{-1/2}t,\qquad
 v'=R^{1/4}(v-v_0),\qquad R'=R^{1/2}.
\]
Then $x-2vt-u_T=R^{1/4}(x'-2v't'-u_T')$ and
$|\det\Phi|=R^{-(k+2)/4}$.  A fine cap and its $\Lambda/R$ normal
thickness become a cap of width $R'^{-1/2}$ and thickness $\Lambda/R'$;
an $aT$ or $bT$ becomes the corresponding dilate at radius $R'$ with the
same $\Lambda$.  The transformed fine labels form an allowed translated
lattice.  The Jacobian occurs on both sides of the desired estimate and
cancels.  A lower-scale $R'^{1/2}$-cube pulls back to a cell $C$ of
transverse width $O_d(R^{1/2})$ and longitudinal length $O(R^{3/4})$.

Apply \eqref{p:bd-local} with $s=R^{1/2}$ on an original
$R^{1/2}$-cube $Q$.  Choose its neighborhood radius
$h=\Lambda^{1/4}R^{1/2}$ and take
$L\geq C(B+d)/\vartheta$.  Outside that neighborhood the decoupling
weight is $O(R^{-B-d-10})$.  The relative packet tails make boxes missing
the neighborhood negligible too.  There are at most
$C_d\Lambda^{d/4}$ neighboring $R^{1/2}$-cubes $Q''$, only $O_d(1)$
parent boxes of a fixed $J$ meet the neighborhood, and each $Q''$ meets
only $O_d(1)$ rescaling cells $C$ in such a box.  Thus, writing
$F_\Box=\sum_{T\mapsto\Box}F_T$,
\begin{equation}\label{p:expanded-decoupling}
 \|F\|_{L^p(Q)}\leq C_dD R^{\delta/2}
 \left(\sum_{(\Box,Q'',C)}
       \|F_\Box\|_{L^p(Q''\cap C)}^2\right)^{1/2}
       +R^{-B/2}\mathcal R,
\end{equation}
where $\mathcal R$ denotes the right side of the defining refined
inequality without its constant.  Finite truncation and the polynomial
number of packets and cubes in the local box justify the displayed
negligible error after increasing $B$.

Put $b'=b-C_d\Lambda^{-3/4}$.  If $b'T$ meets a cell $C$ that meets a
neighbor $Q''$ of $Q$, compare the point of intersection with the center
of $Q$ in the adapted coordinates $y=x-2v_0t-u$.  The change in $y$ is
$O_d(\Lambda^{1/4}R^{1/2})$; the change in $t$ is
$O_d(R^{3/4}+\Lambda^{1/4}R^{1/2})$.  Since
$|v-v_0|=O_d(R^{-1/4})$, the transverse residual for the fine tube
changes by at most $C_d\Lambda^{1/4}R^{1/2}$.  Dividing by its width
$\Lambda R^{1/2}$ gives $C_d\Lambda^{-3/4}$.  All strips use the same
longitudinal slab, so this implies $bT$ meets $Q$; no individual tube
endpoint is crossed.

For completeness, here is the counting step.  Pigeonhole the cube norms,
parent packet counts $W_\Box$, lower-scale cell counts
$m_{\Box,b'}(C)\asymp M'$, and numbers $n(Q)\asymp n$ of retained triples.
The polynomial configuration size makes the cost $(\log R)^{C_d}$.
By \eqref{p:expanded-decoupling} and
$(\sum_{i=1}^n a_i^2)^{p/2}\leq n^{p/2-1}\sum_i a_i^p$, a retained class
satisfies
\begin{equation}\label{p:refined-recurrence-a}
 \|F\|_{L^p(U)}^p\lesssim
 (\log R)^{C_d}(DR^{\delta/2})^p
 n^{p/2-1}\Lambda^{d/4}
 \sum_\Box\|F_\Box\|_{L^p(U_{\Box,M'})}^p
 +\text{negligible}.
\end{equation}
A fixed triple is associated with at most $C_d\Lambda^{d/4}$ original
test cubes.  Rescaling and applying $\mathfrak A_d(R',\Lambda,a,b')$
gives, since the original packet norms are comparable and the assigned
families are disjoint,
\[
 \sum_\Box\|F_\Box\|_{L^p(U_{\Box,M'})}^p
 \lesssim\mathfrak A_d(R',\Lambda,a,b')^p
 \left(\frac{M'}W\right)^{p/2-1}
 \left(\sum_T\|F_T\|_p^2\right)^{p/2},\quad W=\#\{T\}.
\]
For a fixed $Q$, one parent box contributes at most
$C_d\Lambda^{d/4}$ triples.  Choose one retained triple from each of at
least $c_dn\Lambda^{-d/4}$ distinct boxes.  Each contributes $\asymp M'$
different fine packets whose $bT$ meets $Q$ by the preceding geometric
transfer.  Disjoint assignment to boxes therefore gives
$nM'\lesssim\Lambda^{d/4}M$.  Substitution into
\eqref{p:refined-recurrence-a}, followed by a $p$th root, proves
\begin{equation}\label{p:refined-recurrence}
 \mathfrak A_d(R,\Lambda,a,b)
 \lesssim D(\log R)^{C_d}R^{\delta/2}\Lambda^{d/8}
 \mathfrak A_d(R^{1/2},\Lambda,a,b-C_d\Lambda^{-3/4})+1.
\end{equation}
The two neighborhood counts give
$\Lambda^{d/(4p)}\Lambda^{(d/4)(1/2-1/p)}=\Lambda^{d/8}$;
there is no dimension-dependent power of $R$.

Stop at the first $R_j=R^{2^{-j}}$ with
$\Lambda\gtrsim R_j^{1/4}$.  Bounded \emph{unclipped} overlap per label
and $O_d(R_j^{k/2})$ labels give the direct estimate
$\mathfrak A_d(R_j,\Lambda,a,b_j)\lesssim R_j^{k/4}
\lesssim\Lambda^k$.  There are $j=O(\log(1/\vartheta))$ steps.  For
$R$ sufficiently large, their total count-dilation drift
$C_dj\Lambda^{-3/4}$ is less than $1/4$, so the initial margin
$b-a\geq1/2$ remains positive.  Iterating
\eqref{p:refined-recurrence} yields
\[
 \mathfrak A_d(R,\Lambda,a,b)
 \lesssim_{d,p,\delta,\vartheta}
 (\log R)^{C_dj}R^{C_d\delta}
 \Lambda^{k+dj/8}.
\]
Choose $\vartheta$ first so that
$\vartheta O_d(\log(1/\vartheta))<\varepsilon/3$, then $\delta$ so that
the decoupling factors cost less than $R^{\varepsilon/3}$, and absorb the
logarithms into the remaining loss.  The weight order $L$ and its fixed
constant are chosen after $\vartheta$; no original packet decay order is
being chosen as a function of $\varepsilon$.  This proves
\eqref{p:buffered-equation}.
\end{proof}

\begin{proposition}[Arbitrary auxiliary weights]\label{p:weighted}
Let $2<p\leq q_k$, and let $f_T$ be band-limited packets of scale $N$.
For $b_T\geq0$, require $f_T=0$ if $b_T=0$, and set
$B_b(z)^2=\sum_Tb_T^2w_T(z)$.  Then, for every $\eps>0$,
\begin{equation}\label{p:weighted-equation}
 \int B_b(z)^{2-p}\left|\sum_Tf_T(z)\right|^p\dd z
 \leq C_\eps N^\eps N^{k+2}\sum_Ts_T^p b_T^{2-p}.
\end{equation}
The integrand is defined as zero where both the field and $B_b$ vanish.
The estimate admits unit suprema.  If all frequency labels lie within
$C/N$ of a $\kappa$-dimensional affine plane, $1\leq\kappa\leq k$, the
range extends to $2<p\leq q_\kappa$.  In the zero-dimensional case there
are only boundedly many frequency slots and the assertion holds for every
fixed finite $p>2$.
\end{proposition}
\begin{proof}
We derive the weighted conclusion from \Cref{p:buffered-refined}, applying
that lemma to relative-tail atoms, not to the unsplit polynomial-tail
$f_T$ with a single original-tube count.  Fix $\eps>0$, put $R=N^2$, and
choose $\vartheta>0$ small enough, after the fixed packet exponents, to
pay every power of $\Lambda=R^\vartheta$ below.  The original profile has
fixed native Fourier support.  Choose a reciprocal sampling lattice with
fixed spacing $\gamma$ in the $k$ native spatial variables and spacing
$\gamma/\Lambda$ in native time.  Smooth Fourier cutoffs equal to one on
the original Fourier box and supported inside the reciprocal fundamental
domains give, by Poisson summation, fixed Schwartz kernels $K_x,K_t$ and
the exact expansion
\begin{equation}\label{p:anisotropic-sampling}
 F_T(X)=\sum_{j\in\Z^k,\,m\in\Z}c_{T,j,m}
 K_x(X_x-u_{T,x}-\gamma j)
 K_t\bigl(\Lambda(X_t-u_{T,t})-\gamma m\bigr),
\end{equation}
where $u_T$ is the native cell center and, with a fixed normalization,
\[
 c_{T,j,m}=C_\gamma F_T(u_{T,x}+\gamma j,
                              u_{T,t}+\gamma m/\Lambda),\qquad
 |c_{T,j,m}|\leq C s_T
 (1+|j|+|m|/\Lambda)^{-P'}.
\]
The temporal Fourier support of an atom is $O(\Lambda)$ in native
coordinates.  In physical coordinates the Fourier map is
\[
 \xi=v+\eta/N,\quad
 \tau=-|v|^2-2v\cdot\eta/N+\sigma/N^2,\quad
 \tau+|\xi|^2=(\sigma+|\eta|^2)/N^2.
\]
Thus an atom has tangential cap width $O(R^{-1/2})$ and normal thickness
$O(\Lambda/R)$, exactly the range of \Cref{p:buffered-refined}.  Its
$L^p$ norm is comparable to
$|c_{T,j,m}|N^{(k+2)/p}\Lambda^{-1/p}$, with no lower bound on
$|c_{T,j,m}|$ required.  A fixed-kernel sampling in time would not give
the required relative $R^{-B}$ tail outside an $R$-long tube; the
$\Lambda^{-1}$ native-time compression in
\eqref{p:anisotropic-sampling} does.

Fix one sample index $(j,m)$ and write $g_T$ for its atom.  Let
$d_{T,j,m}$ be native distance to its shifted center in the
\emph{original} coordinates $X_{N,v}$ of \eqref{p:native}, and set
\[
 w_{T,j,m}=(1+d_{T,j,m})^{-P},\qquad
 B_{b,j,m}^2=\sum_Tb_T^2w_{T,j,m}.
\]
Partition spacetime into test boxes with time cores of length $R$ and
bounded spatial overlap.  For a core centered at $t_0$, retain atoms whose sample
time centers lie in $[t_0-2R,t_0+2R]$, and use the \emph{common} support
slab $S=[t_0-5R,t_0+5R]$ for all retained atoms.  Test cubes lie in
$[t_0-R/2,t_0+R/2]$.  Thus retained centers have at least a $3R$
margin to $\partial S$, while omitted centers have at least a fixed
$R$ separation from the test cubes.  The retained atom has
$R^{-B}$ relative time tail outside $S$; the omitted atoms gain
$C_A\Lambda^{-A}$ on the core, with a summable further factor in their
time-slab distance.  Retain spatially only trajectories that can meet a
fixed padded box; the other atoms have an arbitrarily small Schwartz
spatial tail.  The phase-space count \eqref{p:count} sums both discarded
classes: with the full shifted envelope kept in the denominator,
three-factor H\"older as below and annular counting give an arbitrary
$\Lambda^{-A_1}$ or $N^{-A_1}$ gain in their weighted $L^p$ norms,
respectively.  This charge is made before applying the local lemma.
A retained atom is concentrated, relative to its own norm,
on the common-slab strip of transverse radius $\Lambda N$, with
$C_A\Lambda^{-A}$ tail outside a fixed support dilate $aT$, say $a=2$.
Count a larger dilate, say $bT$ with $b=3$.  Since $A$ may be chosen after
$\vartheta$, this is $R^{-B}$ for the $B$ of
\Cref{p:buffered-refined}.  Each atom belongs to only
$\Lambda^{O_d(1)}$ padded test boxes.

The actual $v$ may move within its $N^{-1}$ frequency slot.  Relabel it
by the slot center and use one of a fixed number of cap-grid colors.
Over an $O(R)$ slab, the two trajectories differ by $O(N)$, which is
$O(\Lambda^{-1})$ of the inflated transverse radius; the fixed
$a=2$, $b=3$ collar absorbs this displacement.  For this fixed sample
index, the shifted phase-space multiplicity is bounded by a polynomial
in $J=1+|j|+|m|/\Lambda$.  Color intercepts on the \emph{unclipped}
strips into at most $C\Lambda^{C_d}J^{C_d}$ families of bounded
same-label overlap.  There are only $O(1)$ relevant time slabs per test
cube for each such family.  In every local box there are polynomially
many packets.  Set $a_T=|c_{T,j,m}|b_T^{2/p-1}$ for $b_T>0$ and
$a_* = \max_T a_T$ in that box.  Discard only the classes with
$a_T<N^{-C_0}a_*$, where $C_0$ is chosen after the polynomial local
packet count.  This is a cutoff in the \emph{weighted} amplitude, not in
$|c_{T,j,m}|$.  Indeed, writing $h_T=|g_T|/|c_{T,j,m}|$ and taking a
Schwartz majorant $h_T\lesssim w_{T,j,m}$, three-factor H\"older gives
\[
 B_{b,j,m}^{2/p-1}\sum_{a_T<N^{-C_0}a_*}|g_T|
 \lesssim
 \left(\sum_{a_T<N^{-C_0}a_*}a_T^ph_T\right)^{1/p}
 \left(\sum_T h_T\right)^{1/2}.
\]
The last sum and the number of local packets are polynomially bounded;
$\int h_T\lesssim N^{k+2}/\Lambda$.  Thus this discarded contribution
is $N^{-C_0+O_d(1)}\Lambda^{O_d(1)}J^{O_d(1)}$ times
$(N^{k+2}/\Lambda)^{1/p}(\sum_Ta_T^p)^{1/p}$.
The retained $a_T$ occupy $O(\log N)$ dyadic classes.  Within each
$b_T\asymp\alpha$ class, comparability of $a_T$ also makes the atom
$L^p$ norms comparable, so \Cref{p:buffered-refined} applies.  No
comparison between the multiplicity of shifted atom tubes and the
multiplicity of the original tubes is used.

To see the weighted reduction quantitatively, sort $N$-cubes by
$B_{b,j,m}\asymp\beta$ and atoms by
$b_T\asymp\alpha$, then by retained $a_T$.  The slot-center relabeling
above changes only the local counted strips; the shifted envelope here
keeps $v$ and therefore compares globally with $B_b$.  No absolute cutoff on
$\alpha$ or $\beta$ is needed: the cube classes are disjoint, and for
each $T$ the moderate relative range below contains only $O(\log N)$
$\beta$-classes.  On a cube meeting a counted $bT$, its atom center is
within $O(\Lambda)$ in native space; hence the atom contributes at least
$c\alpha^2\Lambda^{-P}$ to $B_{b,j,m}^2$.  The count $M_{j,m}$ for that
class obeys
\[
 M_{j,m}\lesssim\Lambda^P\beta^2/\alpha^2,\qquad
 \beta^{2-p}M_{j,m}^{p/2-1}
 \lesssim\Lambda^{P(p/2-1)}\alpha^{2-p}.
\]
Use \Cref{p:buffered-refined} on the moderate ratios
$N^{-C}\leq\alpha/\beta\leq C\Lambda^{P/2}$ and sum their
$O(\log N)$ relative ranges.  For smaller ratios the identity
\begin{equation}\label{p:small-ratio}
 B_{b,j,m}^{2/p-1}|c_{T,j,m}|
 =|c_{T,j,m}|b_T^{2/p-1}
       (b_T/B_{b,j,m})^{1-2/p}
\end{equation}
gives $N^{-C(1-2/p)}$; choose $C$ to dominate the polynomial local
packet count and use pointwise H\"older.  For the complement of the
counted activity strips, retain the \emph{anisotropic} atom envelope
\[
 h_{A,\Lambda}(z)
 =(1+|X_x-u_{T,x}-\gamma j|)^{-A}
  (1+\Lambda|X_t-u_{T,t}-\gamma m/\Lambda|)^{-A}.
\]
It gains $\Lambda^{-A_1}$ outside the strip in either transverse or
temporal direction, while $\int h_{A,\Lambda}^p\lesssim
N^{k+2}/\Lambda$.  By \eqref{p:count}, annular counting, and pointwise
H\"older, the off-activity $p$th power is bounded by
$C_A N^{C_p}J^{C_p}\Lambda^{-A_1}(N^{k+2}/\Lambda)
\sum_T|c_{T,j,m}|^pb_T^{2-p}$.  Choose $A$ after $\vartheta$ to absorb
the polynomial factor.  Keeping the temporal factor in
$h_{A,\Lambda}$ is essential to retain $\Lambda^{-1}$ in this bound.
The fixed-index conclusion, including local-box and color summation, is
\begin{equation}\label{p:weighted-atom}
 \int B_{b,j,m}^{2-p}\left|\sum_Tg_T\right|^p
 \leq C_\eps N^{\eps/2}\Lambda^{C(d,p,P)}J^{C(d,p,P)}
 \frac{N^{k+2}}\Lambda
 \sum_T|c_{T,j,m}|^pb_T^{2-p}.
\end{equation}
Here and below the exponents $C(d,p,P)$ are fixed before $\vartheta$.

Native-distance comparison gives
$B_{b,j,m}\leq CJ^{P/2}B_b$.  Since $2-p<0$,
$B_b^{2-p}\leq CJ^{P(p-2)/2}B_{b,j,m}^{2-p}$.  Apply
Minkowski in the original weighted $L^p$ norm and use
\eqref{p:weighted-atom} and the sampled coefficient bound.  The
$m$-sum of $\Lambda^{-1/p}(1+|m|/\Lambda)^{-P'+C}$ is
$O(\Lambda^{1-1/p})$; the $j$-sum converges if the original fixed
$P'$ is chosen greater than $C(d,p,P)+k+1$.  Consequently
\[
 \left\|B_b^{2/p-1}\sum_Tf_T\right\|_p
 \leq C_\eps N^{\eps/(2p)}
 \Lambda^{C(d,p,P)+1-1/p}N^{(k+2)/p}
 \left(\sum_Ts_T^pb_T^{2-p}\right)^{1/p}.
\]
Choose $\vartheta$ sufficiently small that the displayed $\Lambda$
power, all preceding color and logarithmic losses, and the deliberately
smaller buffered loss fit below $N^{\eps/p}$.  Taking the $p$th power
proves \eqref{p:weighted-equation}.  The original $P'$ is chosen once,
independently of $\eps$; only the output Schwartz order and weighted
decoupling order depend on the chosen $\vartheta$.  Unit suprema follow
from \Cref{p:reproduction} and fixed-scale envelope comparability.

Finally rotate an affine $\kappa$-plane to the first $\kappa$ coordinates
and remove its constant and linear normal terms by modulation and shear.
The normal frequencies then have size $O(N^{-1})$, and their quadratic
contribution is $O(N^{-2})$.  Put $m=k-\kappa$ and write $x=(u,y)$ in
tangential and normal coordinates.  For a packet $T$, let $(y_T,t_T)$ be
the normal position and time of its physical center and, on a fixed normal
slice $y$, put $j_T(y)=\lfloor(y-y_T)/N\rfloor\in\Z^m$ coordinatewise.
Write $d_{T,\kappa}$ and $w_{T,\kappa}=(1+d_{T,\kappa})^{-P}$ for the
projected tangential-time native cell.  Since the normal frequency is
$O(N^{-1})$, its native displacement at time $t$ is
$(y-y_T)/N-2(Nv_T^\perp)(t-t_T)/N^2$.  Consequently, uniformly in $t$,
\[
 1+d_T\asymp1+d_{T,\kappa}+|j_T(y)|,
 \qquad w_T\gtrsim(1+|j_T(y)|)^{-P}w_{T,\kappa}.
\]
Choose fixed $A,P''$ with $A+P''\le P'$ and $P''$ sufficiently large for
the $\kappa$-dimensional estimate.  On the slice $y$, the packet has
projected amplitude at most $Cs_T(1+|j_T(y)|)^{-A}$ and decay order
$P''$.  Its projected native Fourier box is fixed: substitution in the
normal native variable shifts time frequency by $O(Nv_T^\perp)=O(1)$,
while the normal quadratic carrier shifts it by
$N^2|v_T^\perp|^2=O(1)$.

Group the sliced packets by $j_T(y)=j$, and write $V_j$ for their sum.
For fixed $y,j$ and each $N^2$
time slab, this condition restricts the normal intercept to $O(1)$
cells for every projected frequency/intercept cell; there are also only
$O(1)$ normal frequency slots.  Hence each $j$-family has bounded
projected phase-space multiplicity.  Its tangential frequency slots
are split into $O(1)$ colors and snapped to a $\kappa$-dimensional
$N^{-1}$ grid.  Use the dimension-$\kappa$ estimate already obtained
from the preceding full-dimensional argument (or induct on $k$).
If
$B_j^2=\sum_{j_T(y)=j}b_T^2w_{T,\kappa}$, then
$B_b^2\gtrsim(1+|j|)^{-P}B_j^2$.  As $2-p<0$, applying the
$\kappa$-dimensional estimate to that family gives
\[
 \int_{u,t} B_b^{2-p}|V_j|^p
 \lesssim N^{\eps+\kappa+2}
 (1+|j|)^{-Ap+P(p-2)/2}
 \sum_{j_T(y)=j}s_T^pb_T^{2-p}.
\]
For each $T,j$, the set of $y$ with $j_T(y)=j$ has volume $N^m$.
Integrating in $y$ supplies exactly $N^m$, and Minkowski sums the
$j\in\Z^m$ contributions provided $A>m+P(p-2)/(2p)$.  This is possible
with the fixed sufficiently large original $P'$, and yields
\eqref{p:weighted-equation} in the range $2<p\le q_\kappa$.
If $\kappa=0$, put $\phi_T=(1+d_T)^{-P'}$.  The bounded
number of slots and phase-space counting give $\sum_T\phi_T\leq C$.
Three-factor H\"older gives
\[
 \Big(\sum_Ts_T\phi_T\Big)^p
 \leq\Big(\sum_Ts_T^pb_T^{2-p}\phi_T\Big)
       \Big(\sum_Tb_T^2\phi_T\Big)^{(p-2)/2}
       \Big(\sum_T\phi_T\Big)^{p/2}.
\]
Since $\phi_T\leq w_T$, multiplication by $B_b^{2-p}$ and integration,
using $\int\phi_T\leq CN^{k+2}$, prove the assertion for every finite $p>2$.
\end{proof}

\subsection{Local estimates without a scale loss}

\begin{proposition}[Local linear estimate]\label{p:local-linear}
Let $V=\sum_Tf_T$ be either a band-limited or a smooth packet field of scale
$N$, and let $B^2=\sum_Ts_T^2w_T$.  For a ball $D$ of radius $N$ and any
$z_0\in D$,
\begin{equation}\label{p:local-linear-equation}
 \|V\|_{L^{q_k}(D)}\leq C N^{k/2}B(z_0).
\end{equation}
The same bound holds for the unit-supremum norm on $D$.
\end{proposition}

\begin{proposition}[Local multilinear estimate]\label{p:local-multilinear}
Let $V_j$ and $B_j$ be as in \Cref{p:local-linear}, with frequency labels
in sets satisfying \eqref{p:transversality}.  For fixed $p>p_k$, a ball $D$
of radius $N$, and arbitrary points $z_j\in D$,
\begin{equation}\label{p:local-multilinear-equation}
 \left\|\prod_{j=1}^{k+1}|V_j|^{1/(k+1)}\right\|_{L^p(D)}
 \leq C_p\delta^{-C_p}N^{k/2}
 \prod_{j=1}^{k+1}B_j(z_j)^{1/(k+1)}.
\end{equation}
Separate unit suprema may be taken for the $k+1$ factors.  In particular the
constant has no $N^\eps$ factor.
\end{proposition}
\begin{proof}[Proof of \Cref{p:local-linear,p:local-multilinear}]
Write
$Eg(x,t)=\int e^{i(x\cdot\xi-t|\xi|^2)}g(\xi)\dd\xi$.
We reduce the packets on $D$ to exact extensions, keeping track of the norm.
Translate the center of $D$ to $(x_0,t_0)$.  For a frequency label $v$, take
an integral-normalized smooth bump $g_v$ of width $c/N$, centered at $v$,
and multiply it by $e^{-i(x_0\cdot\xi-t_0|\xi|^2)}$.
On a fixed enlargement of $D$ its extension is its carrier times
\[
 H_v(x,t)=\int\varphi(\eta)
 e^{i(c/N)(x-x_0-2v(t-t_0))\cdot\eta
      -i(c/N)^2(t-t_0)|\eta|^2}\dd\eta.
\]
For fixed sufficiently small $c$, this differs from $1$ by less than $1/2$.
Both $H_v$ and $H_v^{-1}$ have bounded derivatives in the variables
$(x-x_0,t-t_0)/N$.  Also $\|g_v\|_2\leq C N^{k/2}$.
In the multilinear case $c$ can be shrunk by a fixed power of $\delta$ so
that enlarged supports retain determinant at least $\delta/2$; all resulting
constants are polynomial in $\delta^{-1}$.

For a packet $T$ with label $v$, divide its demodulated profile by $H_v$ on
this enlarged ball, multiply by a fixed cutoff, and expand in a Fourier
series in $(x-x_0,t-t_0)/N$.  The resulting identity on $D$ has the form
\[
 f_T(z)=\sum_{m\in\Z^{k+1}}c_{T,m}
 e^{i\ell_m\cdot(z-z_0)/N}Eg_v(z),\qquad
 |c_{T,m}|\leq C_Ls_T(1+d_T(z_0))^{-P''}(1+|m|)^{-L},
\]
where $P''>P/2+k+1$ can be fixed below the available original decay margin.
Only finitely many derivatives are needed for a chosen $L$.  For each mode,
combine the coefficients with the same frequency slot; Cauchy--Schwarz in
the spatial and time labels and their summable native tails gives
\[
 \left\|\sum_Tc_{T,m}g_{v(T)}\right\|_2^2
 \leq C_LN^k(1+|m|)^{-2L}
 \sum_Ts_T^2w_T(z_0).
\]
Here one first uses bounded overlap of the bump supports in distinct
frequency slots; within a slot the native-center lattice has a summable
weight.  This explains why the square bound has no factor equal to the
number of spatial labels.

The exact paraboloid Strichartz inequality
$\|Eg\|_{q_k}\leq C\|g\|_2$ follows here directly from the dispersive
kernel and the $TT^*$ argument.  Indeed the Gaussian oscillatory integral
gives $\|e^{it\Delta}\|_{1\to\infty}\leq C|t|^{-k/2}$, while Plancherel
gives its $L^2$ norm equal to one.  Interpolation at spatial exponent $q_k$
gives decay $|t|^{-k(1/2-1/q_k)}=|t|^{-2/q_k}$.
The one-dimensional Hardy--Littlewood--Sobolev inequality in time consequently
bounds $TT^*:L^{q_k'}(\R^{k+1})\to L^{q_k}(\R^{k+1})$, and duality gives
the stated extension estimate.  Applying it to each mode gives
\eqref{p:local-linear-equation}, after summing the modes.
For the multilinear assertion use \citet[Theorem~1.7]{Tao2020}.
In ambient dimension $d=k+1$, its product exponent $2a$, $a>1/(d-1)$,
becomes geometric-mean exponent $p=2da>p_k$.
Its quantitative constant is polynomial in a parameter controlling domain
size, finitely many graph derivatives, inverse determinant and inverse
interior support distance.  On bounded paraboloid patches these quantities
are bounded by a fixed power of $\delta^{-1}$: enlarge each relevant domain
by $c\delta$, and, if necessary, subdivide into patches of polynomially small
width in $\delta$.  The Hessian is constant, all required higher derivatives
are bounded, and the bump supports lie a polynomial distance inside the
domains.  Thus the theorem applies globally with the displayed
$C_p\delta^{-C_p}$ and no scale loss.
The external modulation belonging to a fixed Fourier mode has modulus one.
Expand the product over modes using subadditivity of the $1/(k+1)$ power,
and then Minkowski when $p\geq1$.  Taking $L$ larger than the dimension
times $k+1$ makes this sum converge.

Finally apply \eqref{p:sup-convolution} separately to the exact extension
factors.  A spacetime translation of $Eg$ is another extension with input
multiplied by a unimodular phase, so it has the same input norm and frequency
support.  Tonelli and the global estimates bound each translated product.
The convolution exponents can be chosen to sum all mode and translation
weights.  This proves the separate-supremum assertions and permits the
independent choices of $z_j$.
\end{proof}

\subsection{Hierarchies, masks, and positive weights}

\begin{lemma}[Sampling preserves the envelope exponent]\label{p:convolution}
Let $K'\geq K\geq1$, let a $K'$-cap lie in a $K$-cap, and let $T'$ be a
child cell of scale $K'$.  If $\mathcal G_K$ is the parent grid in that cap,
then, for $P>k+2$ sufficiently large,
\begin{equation}\label{p:convolution-equation}
 \sum_{T\in\mathcal G_K}w_T(z)w_{T'}(z_T)\leq C_Pw_{T'}(z).
\end{equation}
The same exponent $P$ occurs on both sides, independently of $K'/K$.
If the sum is restricted to $d_T(z)>L$, $L\geq2$, it is at most
\begin{equation}\label{p:convolution-tail}
 C_P\left(L^{-(P-k-1)}+(K/K')^{P-k-2}\right)w_{T'}(z).
\end{equation}
When $d_{T'}(z)$ is bounded, the second term can be omitted.
\end{lemma}
\begin{proof}
Put $r=K/K'$.  In parent coordinates the linear map to child coordinates
has spatial coefficients $O(r)$ and time coefficient $r^2$, determinant
$r^{k+2}$; the possible velocity shear has bounded normalized coefficients
because the child frequency lies in the parent cap.  Let
$a=1+d_{T'}(z)$.  Parent centers whose child distance is at least $a/2$
contribute at most
$Ca^{-P}\sum_Tw_T(z)\leq Ca^{-P}$.  With $d_T(z)>L$, this sum has the
additional lattice-tail factor $L^{-(P-k-1)}$.
For the other centers, necessarily $a\gg1$ and $d_T(z)\geq ca/r$.
Their contribution is at most
\[
 C(a/r)^{-P}\sum_{T\in\mathcal G_K}w_{T'}(z_T)
 \leq Ca^{-P}r^P r^{-(k+2)}.
\]
The last lattice sum is bounded by its integral at the child scale plus a
fixed boundary term, with volume ratio $r^{-(k+2)}$.
This proves both conclusions and shows explicitly why no envelope power is
lost in repeated convolutions.
\end{proof}

\begin{proposition}[Weighted transport through a hierarchy]\label{p:hierarchy}
Take nested frequency grids at scales $K<K'$, rounded within fixed factors.
Assemble all child fields in each parent cap, and localize them on its parent
physical grid.  The localizers can be chosen band-limited, with summable
Schwartz tails and bounded sums, so that the resulting parent packets have
amplitudes $a_T$ and satisfy
\begin{equation}\label{p:parent-weights}
 b_T^2\geq c\sum_{T'\text{ in the parent cap}}b_{T'}^2w_{T'}(z_T)
\end{equation}
whenever the parent weights are chosen to dominate the displayed sum.
For $2<p\leq q_k$,
\begin{equation}\label{p:hierarchy-equation}
 \sum_T K^{k+2}a_T^p b_T^{2-p}
 \leq C_\eps(K'/K)^\eps
 \sum_{T'}(K')^{k+2}a_{T'}^p b_{T'}^{2-p}.
\end{equation}
It remains valid after retaining any subset of parent cells, at every step
of a finite hierarchy.  After modifying a child field, use its fresh
amplitudes in the same fixed parent masks.  On subtrees lying in a
codimension-one strip at their final frequency resolution one may use
$p=q_{k-1}$ at every transition; for $k=1$ any fixed finite $p>2$ is
available.  Products of the transition estimates give the complete
finite-hierarchy estimate.

Fixed polynomial weights on macroscopic boxes, and any fixed additional
edge decay, are allowed.  Their orders are chosen once, before the finite
grid is chosen.  Constants may depend on the finite grid; the required
original profile regularity and decay order do not grow with its depth.
\end{proposition}
\begin{proof}
A concrete band-limited partition is obtained by taking $\chi=|\phi|^2$,
where $\widehat\phi$ is smooth with sufficiently small support.  Poisson
summation makes $\sum_{j\in\Z^{k+1}}\chi(X-j)$ a positive constant after
normalization; alternatively a bounded lattice average gives the same
identity.  These nonnegative localizers have arbitrarily rapid tails,
compact Fourier support, and bounded sums.  If $U$ is the pre-cutoff field
in a parent cap, choose
\begin{equation}\label{p:parent-amplitude}
 a_T=C\sup_z |U(z)|(1+d_T(z))^{-C_0},
\end{equation}
where $C_0$ exceeds the required output decay order by a fixed margin.
Multiplying by the localizer makes the resulting normalized packet satisfy
any prescribed fixed decay order.  Its native Fourier support stays in a
fixed enlarged box.

In parent coordinates the child packets have scale $K'/K$.  Apply
\Cref{p:weighted} there, and then its unit-supremum version to
\eqref{p:parent-amplitude}.  The child square sum varies by at most a fixed
polynomial over parent distance, since velocity differences are $O(K^{-1})$.
The sufficiently large $C_0$ absorbs this variation, while
\eqref{p:parent-weights}, whose exponent $2-p$ is negative, gives the required
upper bound for the parent weighted norm.  Changing coordinates supplies
$K^{k+2}$ and yields \eqref{p:hierarchy-equation}.
Deleting parent outputs decreases its left side.  For modified inputs the
same proof uses the same localizers and recomputes \eqref{p:parent-amplitude};
it never assumes that cancellation or an old amplitude survives deletion.
The affine part of \Cref{p:weighted} gives the larger exponent on strip
subtrees: a strip thin enough at the final precision is thin enough at each
intervening child precision.

We explain the two uniformity assertions needed in iteration.  The new
square background formed by sampling child squares satisfies
\[
 \sum_Tw_T(z)\sum_{T'}b_{T'}^2w_{T'}(z_T)
 \leq C\sum_{T'}b_{T'}^2w_{T'}(z)
\]
by \Cref{p:convolution}; the outer exponent has not decreased.
Additional decay on an edge only improves this bound.  A large displacement
along a finite path forces a large displacement on an edge after normalization
to the largest relevant cell.  Keeping a fixed reserve of decay on every edge
therefore bounds paths with a large displacement by the full-background
convolution times a negative power of that displacement.  The choice of
which edge is large costs at most the number of edges, rather than a new
decay order at each edge.

For a positive macro weight $\Psi$, suppose both ratios of its values at
an edge's endpoints are bounded by $C(1+d)^J$, where
$d=d_{T'}(z_T)$ is the child distance.  The constants and $J$ must be
independent of the main scale; this holds for polynomial weights on macro
boxes whose size dominates the affected cells.  Conjugate the given
denominators by
\[
 \widetilde b_T=b_T\Psi(z_T)^{-1/(p-2)},\qquad
 Q=P+\frac{2J}{p-2}.
\]
Writing $w_{T'}^{(Q)}=(1+d_{T'})^{-Q}$, one has
\begin{align*}
 \sum_{T'}\widetilde b_{T'}^2w_{T'}^{(Q)}(z_T)
 &\leq C\Psi(z_T)^{-2/(p-2)}
             \sum_{T'}b_{T'}^2w_{T'}^{(P)}(z_T)\\
 &\leq C\widetilde b_T^2.
\end{align*}
Apply the one-step proof with this stronger fixed child envelope exponent
and its fixed supremum-decay reserve.  Since
$\widetilde b_T^{2-p}=b_T^{2-p}\Psi(z_T)$, it proves the explicit estimate
\[
 \sum_T\Psi(z_T)K^{k+2}a_T^pb_T^{2-p}
 \leq C_\eps(K'/K)^\eps
 \sum_{T'}\Psi(z_{T'})(K')^{k+2}a_{T'}^pb_{T'}^{2-p}.
\]
The original denominators have been retained.  The same value of $Q$ works
at each edge; the reserve does not accumulate with depth.  For a partition
of macro space with summable polynomial weights, summing over macros costs
their bounded total overlap.  This gives the version for restricted subtrees
and frozen masks, without using discontinuous physical cutoffs.

Finally, for separated nested scales, inverse widths added by localization
form a geometric series, including the bounded velocity shears.  Thus the
fixed Fourier enlargement can be chosen once.  A block of comparable scales
is treated as one transition; finitely many extra masks change constants.
Repeated constants can grow with the fixed grid, but not with the main
scale.  In a later slow diagonal those constants may be absorbed only after
the grid is fixed.  Arbitrarily high Schwartz output decay is supplied by
\Cref{p:reproduction}; it imposes no increasing regularity condition on the
original packets.
\end{proof}

\subsection{A near-Cauchy estimate}

The following estimate has no loss depending on $N$.  Its threshold is
measured against the pointwise Cauchy bound $|V|\leq CN^{k/2}B$; the allowed
loss depends instead on the distance $R$ below that bound.

\begin{theorem}[Near-Cauchy large values]\label{p:near-cauchy}
Let $V=\sum_Tf_T$ be a band-limited field of scale $N$, and let
$B^2=\sum_Ts_T^2w_T$.  Suppose a finite set $Y$ of unit cubes has witnesses
$z_Q\in Q$ such that
\[
 B(z_Q)\leq1,\qquad |V(z_Q)|\geq N^{k/2}/R,\qquad R\geq1.
\]
Then, for every $\eps>0$,
\begin{equation}\label{p:near-cauchy-equation}
 |Y|\leq C_\eps R^{q_k+\eps}\sum_Ts_T^2.
\end{equation}
The conclusion holds with unit suprema and fixed enlargements of the witness
cubes.  The Fourier enlargement and original decay orders can be fixed
before $\eps$ is chosen.
\end{theorem}

We first provide a finite exponent from which the improvement can start.
This also makes the exponent-infimum argument below noncircular.

\begin{lemma}[A polynomial starting bound]\label{p:polynomial-start}
Under the hypotheses of \Cref{p:near-cauchy}, there is a fixed finite $M_k$
such that
\begin{equation}\label{p:polynomial-start-equation}
 |Y|\leq C(2+R)^{M_k}\sum_Ts_T^2.
\end{equation}
\end{lemma}
\begin{proof}
Write $S=\sum_Ts_T^2$.  By \Cref{p:weighted} with $b_T=s_T$ and $p=q_k$,
\begin{equation}\label{p:rough-scale-bound}
 |Y|\leq C_\eps N^\eps R^{q_k}S.
\end{equation}
At a witness, $B\leq1$ and $2-q_k<0$; the unit-supremum version of the
weighted estimate supplies the inequality for arbitrary witness positions.
Consequently any range $N\leq(2+R)^D$, with $D$ fixed, already satisfies
\eqref{p:polynomial-start-equation}.  We may assume $N$ exceeds a sufficiently
large fixed power of $2+R$.

Discard the packets with $s_T>C R N^{-k/2}$.  Their total field at a witness
has modulus at most
\[
 \sum_{s_T>C R N^{-k/2}}s_T(1+d_T)^{-P'}
 \leq (C R N^{-k/2})^{-1}\sum_Ts_T^2w_T
 \leq C^{-1}N^{k/2}/R.
\]
Choose $C$ large, so the remaining field retains half the threshold.
Its coefficients have maximum $s_0=C R N^{-k/2}$.
Put $P_1=(C(2+R))^{10}$, partition the frequency patch into bins of side
$P_1/N$, and write the remaining field as $\sum_bV_b$.
There are $O((N/P_1)^k)$ bins.  Cauchy--Schwarz therefore gives
\begin{equation}\label{p:bin-lower}
 \sum_b|V_b(z_Q)|^2\geq cP_1^k/R^2.
\end{equation}
In this square sum remove all ordered pairs of frequency labels whose
separation is at most $C'/N$, where $C'$ is a sufficiently large fixed
constant compared with the native Fourier box.  The removed expression is
pointwise bounded by $CB^2$: each frequency slot has only boundedly many
such neighbors, and the sum over the spatial and temporal labels is bounded
by Cauchy--Schwarz with the summable packet envelopes.  It is therefore $O(1)$
at the witnesses.  Let $G=\sum_bG_b$ be the sum of the remaining products.
Choosing the constant in $P_1$ large makes \eqref{p:bin-lower} imply
$|G(z_Q)|\geq1$.

We verify the required spacetime almost orthogonality.  If $v_b$ is the
center of a bin, the Fourier support of $G_b$ is contained in
\begin{equation}\label{p:off-diagonal-support}
 c/N\leq|\zeta|\leq CP_1/N,\qquad
 |\sigma+2v_b\cdot\zeta|\leq CP_1^2/N^2.
\end{equation}
To see this, subtract the two cap Fourier coordinates used in the proof
of \Cref{p:weighted}.  Removing the $C'/N$ pairs guarantees the lower bound
on $|\zeta|$.  The deviation from the tangent plane at $v_b$ is the product
of two $O(P_1/N)$ quantities, together with the $O(N^{-2})$ normal errors.
For a fixed $(\zeta,\sigma)$, the bin centers lie in a slab of thickness
$CP_1^2/(N^2|\zeta|)$.  Counting their grid of mesh $P_1/N$ gives overlap
at most $CN^{k-1}P_1^{C_k}$.  Plancherel yields
\[
 \|G\|_2^2\leq C N^{k-1}P_1^{C_k}\sum_b\|G_b\|_2^2.
\]
Within a bin, the effective pointwise packet multiplicity is $O(P_1^k)$.
If $A_b=\sum_{T\in b}s_T(1+d_T)^{-P'}$, then
$|G_b|\leq A_b^2$, $A_b\leq CP_1^ks_0$, and
$A_b^2\leq CP_1^k\sum_{T\in b}s_T^2(1+d_T)^{-P'}$.
Integrating these inequalities proves
\begin{equation}\label{p:polynomial-L2}
 \|G\|_2^2
 \leq C N^{k-1}P_1^{C_k}(R N^{-k/2})^2E
 =C P_1^{C_k}R^2N^{k+1}S.
\end{equation}
All Fourier frequencies of $G$ have size at most $CP_1/N$ in the $k+1$
spacetime variables.  Cover spacetime by cubes of side
$L=N/P_1^{C_2}\geq1$, choosing $C_2$ fixed and sufficiently large.
The rescaled version of \eqref{p:sup-convolution}, summed over this lattice,
shows that the number of these cubes with $\sup|G|\geq1$ is at most
\[
 C L^{-(k+1)}\|G\|_2^2\leq (2+R)^{C_k}S.
\]
Each such cube is contained in an $N$-ball.  If it contains a witness, use
that witness as the envelope reference in \Cref{p:local-linear}; its
unit-supremum estimate bounds the number of high unit cubes in the ball by
$C R^{q_k}$.  Multiplying the last two bounds proves the lemma.
\end{proof}

\begin{proof}[Proof of \Cref{p:near-cauchy}]
Fix the frequency box, Fourier enlargement, phase-space bound and original
regularity class once and for all.  Call an exponent $p$ admissible if
$|Y|\leq C_pR^pS$ for all configurations in this class.  Fixed changes of
multiplicity, amplitude normalization, or cell enlargement change constants
and do not change the infimum of admissible exponents.  By
\Cref{p:polynomial-start} this infimum $p_*$ is finite.  We prove $p_*\leq q$,
where $q=q_k$.

Suppose instead that $p_*>q$.  There are configurations and parameters
$R\to\infty$ with
\begin{equation}\label{p:extremal-sequence}
 |Y|\geq R^{p_*-o(1)}S.
\end{equation}
Indeed use a failing exponent $p_*-\epsilon_j$ and let
$\epsilon_j\downarrow0$, with $R$ large enough to absorb all constants for
the finitely many estimates already chosen.  Bounded $R$ cannot give such
failure by \Cref{p:polynomial-start}.  Equation \eqref{p:rough-scale-bound},
used with each fixed positive $\eps$, now implies
\begin{equation}\label{p:scale-separation}
 \frac{\log N}{\log R}\longrightarrow\infty.
\end{equation}
All subsequent $R^{o(1)}$ statements have this finite-parameter meaning:
first fix positive exponent tolerances and finitely many tests; then take
$R$ sufficiently large, and finally let the tolerances tend to zero.
In particular an admissible exponent just above $p_*$, rather than the
possibly inadmissible exponent $p_*$ itself, is used in every application
below.

\paragraph{Removing excessively occupied fine packets.}
Fix $C_1$ large.  Delete a fine packet if its $CR$ native enlargement contains
more than $R^{C_1}$ witness points.  Double counting with $B(z_Q)\leq1$ gives
\[
 \sum_{T\text{ deleted}}s_T^2
 \leq C R^{P-C_1}|Y|.
\]
Apply any fixed admissible estimate of exponent $p>p_*$ to the deleted
subfield, with threshold a small fixed multiple of $N^{k/2}/R$; its square
envelope is bounded by the original one.  The number of points where this
subfield reaches that threshold is at most
$C R^{p+P-C_1}|Y|$.  Choose $C_1>p+P+10$.
Except on a vanishing fraction of $Y$, the remaining field still reaches a
fixed multiple of the old threshold.  Constants of this sort are absorbed
by changing $R$ by a fixed factor.  Every surviving fine packet now has
occupancy at most $R^{C_1}$ in its $CR$ enlargement.

\paragraph{Coarse assembly and saturation of a ratio class.}
Choose a fixed $A$ much larger than $C_1$ and all preceding fixed exponents,
and put
\[
 K=N/R^A,\qquad L=N/K=R^A.
\]
Round to nested grids within bounded factors.  Equation
\eqref{p:scale-separation} ensures $K\to\infty$ faster than every fixed
power of $R$.  Assemble the surviving fine fields into $K$-packets by
\Cref{p:hierarchy}.  In the sampling definition of their backgrounds use
one slightly larger, fixed fine weight exponent $P+\sigma$, $\sigma>0$:
\begin{equation}\label{p:coarse-background}
 b_T^2=C\sum_{T'\text{ in its cap}}s_{T'}^2
                   (1+d_{T'}(z_T))^{-P-\sigma}.
\end{equation}
The original decay margin permits this choice.  Apply
\Cref{p:weighted,p:hierarchy} with this exponent for the fine estimate and
with the original exponent $P$ for the outer parent envelope.  By
\Cref{p:convolution},
\begin{equation}\label{p:coarse-budgets}
 \sum_Tb_T^2w_T(z)\leq CB(z)^2,\qquad
 \sum_TK^{k+2}a_T^qb_T^{2-q}\leq R^{o(1)}E.
\end{equation}
The ratio satisfies $a_T/b_T\leq CL^{k/2}$ by pointwise Cauchy--Schwarz
in the $L^k$ fine frequency slots and summability over their native centers.
The global contribution of ratios less than $L^{k/2}R^{-2}$ is at most
$C K^{k/2}L^{k/2}R^{-2}B$, again by Cauchy--Schwarz, so is negligible at
the witnesses.  Only $O(\log R)$ dyadic ratio classes remain.  By the
triangle inequality one class
\[
 a_T\asymp r b_T
\]
retains threshold $N^{k/2}/R^{1+o(1)}$ on at least $R^{-o(1)}|Y|$ points.
Use this class henceforth.  Its amplitude-square envelope is $O(r^2)$, and
\eqref{p:coarse-budgets} bounds its energy by $R^{o(1)}r^{2-q}E$.
After dividing its field by $Cr$, its effective threshold parameter at scale
$K$ is $R^{1+o(1)}rL^{-k/2}$.  This parameter is bounded below by a fixed
constant whenever the threshold is attained, by the pointwise Cauchy bound.
Apply an admissible exponent tending to $p_*$.  Since $r$ and $L$ are fixed
powers of $R$, its small exponent error is absorbed in $R^{o(1)}$, giving
\begin{align}
 |Y|&\leq R^{o(1)}(RrL^{-k/2})^{p_*}
              E K^{-(k+2)}r^{-q}\notag\\
 &=R^{p_*+o(1)}S\,(rL^{-k/2})^{p_*-q}.
 \label{p:ratio-saturation}
\end{align}
Here $kq/2=k+2$.  Comparison with \eqref{p:extremal-sequence}, and the
upper bound $r\leq CL^{k/2}$, force
\begin{equation}\label{p:saturated-ratio}
 r=L^{k/2}R^{o(1)},\qquad
 \sum_{T\text{ selected}}b_T^2
 \leq R^{o(1)}E K^{-(k+2)}r^{-q}
 \leq R^{o(1)}S.
\end{equation}
The same fixed packet class remains available at this coarser scale: the
inherited Fourier enlargement shrinks by $L^{-1}$ in parent coordinates,
whereas the parent localizer adds a fixed amount.  A sufficiently large
initial enlargement is consequently preserved as $L\to\infty$.  Its
Schwartz profiles have the originally fixed required bounds after constant
renormalization.  This verifies the uniformity needed to use the same $p_*$.

\paragraph{Localizing the selected square mass.}
We detail the restrictions used to draw a packet through a point.  First
discard globally all selected labels with
$b_T>R^3K^{-k/2}$.  Since the packet decay exceeds $P$,
\[
 \left|\sum_{b_T>R^3K^{-k/2}}f_T(z)\right|
 \leq C r K^{k/2}R^{-3}\sum_Tb_T^2w_T(z)
 \leq C r K^{k/2}R^{-3}
\]
at the witnesses.  By \eqref{p:saturated-ratio} this is negligible relative
to $N^{k/2}/R$.  This global estimate, made before restricting distances,
avoids any unsupported estimate of field tails from a square bound alone.

The remaining selected square envelope is at least $R^{-o(1)}$ on
$R^{-o(1)}|Y|$ retained points.  To verify this, fix $\eta>0$ and consider
points where that envelope squared is at most $R^{-\eta}$.  Normalize the
field there by $rR^{-\eta/2}$ and apply an admissible near-Cauchy exponent.
Compared with \eqref{p:ratio-saturation}, the threshold parameter gains
$R^{-\eta/2}$ and the coefficient-square sum costs $R^\eta$.
The count of these points is therefore at most
\[
 R^{p_*+o(1)}S\,R^{-\eta(p_*/2-1)}.
\]
Since $p_*>q>2$, this is negligible compared with
\eqref{p:extremal-sequence}.  Letting $\eta$ tend to zero along the finite-test
diagonal proves the asserted lower bound.

For any fixed small $\eta>0$, \Cref{p:convolution} gives, term by term in
the original fine square sum,
\[
 \sum_{T:d_T(z)>R^\eta}b_T^2w_T(z)
 \leq C\big(R^{-\eta(P-k-1)}+R^{-A(P-k-2)}\big)B(z)^2.
\]
For fine distance $a\gg1$, the second term comes only from parent centers
moving to within $a/2$ of the fine core; they require parent distance
$\gtrsim aN/K$ and cost $R^{-A(P-k-2)}a^{-P}$.
For bounded $a$ only the first term occurs.  Thus, choosing the square-mass
tolerance smaller than these fixed gains, we may restrict to
$d_T(z)\leq R^{o(1)}$ and retain square mass $R^{-o(1)}$.

We may also require each such selected $T$ to have a surviving fine packet
of its cap within $R^{o(1)}$ fine native distance of $z$.
Indeed, if all its fine packets have distance greater than $R^{c\eta}$,
then parent centers at distance at most $R^\eta$ from $z$ still have those
fine distances greater than $\tfrac12R^{c\eta}$, since the displacement in
fine coordinates is $O(L^{-1}R^\eta)$.  The extra fixed $\sigma$ in
\eqref{p:coarse-background} supplies $R^{-c\eta\sigma}$ when it is dropped
back to exponent $P$.  Sum against the parent weights and use
\Cref{p:convolution}.  These unattached labels have negligible square mass.
This is why one larger exponent was chosen in the sampling definition; no
stronger square-envelope hypothesis at the original witnesses is needed.

Let $I_0(z,T)$ denote the resulting incidence relation.  With
\[
 m_0(z)=\sum_Tb_T^2I_0(z,T),
\]
we have, on the retained set still called $Y$,
\begin{equation}\label{p:incidence-mass}
 R^{-o(1)}\leq m_0(z)\leq R^{o(1)},\quad
 b_T\leq R^{3+o(1)}K^{-k/2},\quad
 \sum_Tb_T^2\leq R^{o(1)}S.
\end{equation}
The upper bound follows by removing $w_T$ at distance $R^{o(1)}$ in
\eqref{p:coarse-budgets}; the lower bound follows from the localized square
mass.  Draw $T$ at $z$ with probability $b_T^2I_0(z,T)/m_0(z)$ and attach
one of its nearby surviving fine packets.  At a fixed $z$ there are only
$R^{o(1)}$ local parent cells in each angular $1/K$ slot.  Consequently its
angular law $\mu_z$ satisfies
\begin{equation}\label{p:slot-upper}
 \mu_z(\text{one }1/K\text{ slot})\leq R^{6+o(1)}K^{-k}.
\end{equation}

\paragraph{Typical directions and resampling.}
Restrict the draw to directions $v$ satisfying
\begin{equation}\label{p:ball-lower}
 \mu_z(B(v,s))\geq R^{-1-o(1)}s^k
 \qquad(K^{-1}\leq s\leq1).
\end{equation}
This removes $O(R^{-1+o(1)})$ probability, without a logarithm of the number
of scales.  Here is a finite covering proof.  In a fixed finite collection
of shifted dyadic grids, restrict attention to cells meeting the bounded
frequency box, with side between $c/K$ and a fixed constant.  Select the
maximal cells in this finite collection for which
$\mu_z(Q)<cR^{-1}|Q|$.  Within each grid these cells are disjoint, so the
sum of their masses is at most $CR^{-1}$.  A ball failing
\eqref{p:ball-lower} contains a cell of comparable side through its center
in one of these grids; that cell belongs to a maximal bad cell.  The bounded
number of grids and fixed comparison factors are absorbed by the indicated
subpower tolerance.  Renormalizing the restricted law costs a bounded
factor.  Let $I(z,T)$ be the resulting incidence relation and
$m(z)=\sum_Tb_T^2I(z,T)$; \eqref{p:incidence-mass} remains true.

Fix $z\in Y$.  For another point $y\in Y$ write $\tau=|t_y-t_z|$.
If $\tau\geq N$, the probability, under the restricted draw, that $y$ is in
the same $R^{o(1)}$-enlarged $K$-cell as $z$ is at most
\begin{equation}\label{p:coarse-hit}
 R^{6+o(1)}(K/\tau)^k.
\end{equation}
Indeed its directions lie in a ball of radius $KR^{o(1)}/\tau$;
\eqref{p:slot-upper} counts the slots in that ball.  A nonempty incidence
also requires $\tau\leq K^2R^{o(1)}$, so the lowest angular precision
$1/K$ costs only a further $R^{o(1)}$ at the endpoint of this range.

If this probability is nonzero, choose a typical direction $v$ that sees
$y$.  In the original draw, the ball of radius $cN/\tau$ about $v$ has mass
at least
\begin{equation}\label{p:fine-hit}
 R^{-1-o(1)}(N/\tau)^k
\end{equation}
by \eqref{p:ball-lower}.  The radius lies between $1/K$ and a fixed constant:
$\tau\geq N$ gives the upper bound, and
$\tau\leq K^2R^{o(1)}$ together with $N/K=R^A$ gives the lower one.
Every draw in that ball has an attached fine packet which sees both $z$ and
$y$ in its $CR$ enlargement.  To check this last statement, its intercept at
$z$ differs by at most $NR^{o(1)}$; the change of angular direction over
time $\tau$ costs at most $cN$; and the discrepancy between a parent label
and a fine label is $O(1/K)$, costing at most
$\tau/K\leq KR^{o(1)}\ll N$.  The time separation is also much less than
$N^2$, whereas $z$ was already at fine native distance $R^{o(1)}$.
All these errors are smaller than the available $CR$ enlargement.

Let $P_f(y)$ be the probability that the attached fine packet in the original
draw sees $y$ in that enlargement.  Comparing
\eqref{p:coarse-hit} and \eqref{p:fine-hit} bounds the expected number of far
points in the sampled coarse cell by
\begin{equation}\label{p:far-occupancy}
 \sum_{\tau\geq N}R^{7+o(1)}(K/N)^kP_f(y)
 \leq R^{C_1+7-Ak+o(1)}.
\end{equation}
The last inequality uses the fine occupancy cutoff: every attached fine
packet sees at most $R^{C_1}$ points, so $\sum_yP_f(y)\leq R^{C_1}$.
This is a comparison with the original draw, not an independence assertion
about a direction and a point.

For $\tau<N$, a point in the same enlarged coarse cell lies in a ball of
radius $CN$ about $z$, since the velocities are bounded and
$KR^{o(1)}\ll N$.  The local linear estimate at $z$, or a bounded cover by
$N$-balls and envelope comparability, bounds the number of such high points
by $CR^q$.  Choose $A$ so large that $Ak>C_1+q+10$.  Combining the near and
far parts, the expected occupancy of the sampled enlarged coarse cell,
for every retained $z$, is at most
\begin{equation}\label{p:occupancy-upper}
 R^{q+o(1)}.
\end{equation}

\paragraph{The contradictory lower occupancy.}
Put $n_T=\#\{z\in Y:I(z,T)=1\}$.  This is no larger than the geometric
occupancy of the enlarged coarse cell.  Average that occupancy first over
the draw at $z$, then uniformly over $z\in Y$.  By
\eqref{p:incidence-mass} and Cauchy--Schwarz, the result is at least
\begin{align*}
 \frac{R^{-o(1)}}{|Y|}\sum_Tb_T^2n_T^2
 &\geq\frac{R^{-o(1)}}{|Y|}
       \frac{(\sum_Tb_T^2n_T)^2}{\sum_Tb_T^2}\\
 &\geq R^{-o(1)}\frac{|Y|}{S}
 \geq R^{p_*-o(1)}.
\end{align*}
The middle inequality uses
$\sum_Tb_T^2n_T=\sum_{z\in Y}m(z)\geq R^{-o(1)}|Y|$.
All selections retained a subpower fraction, so the last inequality still
follows from \eqref{p:extremal-sequence}.  This contradicts
\eqref{p:occupancy-upper} because $p_*>q$.  Thus $p_*\leq q$, and every
$q+\eps$ is admissible, proving \eqref{p:near-cauchy-equation}.
The unit-supremum and fixed-enlargement versions used in the proof follow
from \Cref{p:reproduction} and bounded changes of the lattice.
\end{proof}
